\section{Conclusion}

The proof follows one pressure-normalized velocity history from its initial
datum to the endpoint of its maximal smooth interval.  Every finite temporal
depth and every finite Sobolev depth determine a whole-terminal rectangle,
and the restriction identities assemble their shared entries without changing
the history.  The rectangle estimates then give
\(u\in L^2(0,T_*;H^m)\) for every finite \(m\).  The equation upgrades the
zeroth row to an \(H^m\)-valued endpoint, and the pressure formula together
with the velocity recurrence recovers the complete endpoint jet.  Local
well-posedness restarts from that state, while uniqueness on the overlap keeps
the restarted solution on the history already generated by the datum.  If
\(T_*\) were finite, this continuation would contradict maximality.  Hence the
maximal interval is \([0,\infty)\).

These rectangles also quantify the critical-height rows, frequency tails,
Lipschitz action, and relative-energy stability, with constants depending on
the datum, viscosity, and horizon.  In the small-critical regime, the
scale-invariant norm closes the critical stock/action estimate, and the higher
rectangles admit explicit formulas using the corresponding finite collection
of higher datum norms.

The resulting global solution is smooth, satisfies the exact energy identity,
and carries the decay-normalized pressure.  Relative energy then identifies
every Leray--Hopf solution with that datum with this smooth velocity, and
the pressure is unique up to the unavoidable time-dependent gauge unless this decay normalization is imposed.  The endpoint contradiction and these
consequences all belong to the one fixed-viscosity evolution constructed from
\(u_0\).
